% Epilogue of the monograph (input by AG_monograph.tex)
\chapter*{Epilogue: what is open}
\addcontentsline{toc}{chapter}{Epilogue: what is open}
\markboth{Epilogue}{Epilogue}

Each problem below is left open by a specific chapter, and each is stated in the terms of that
chapter.

\paragraph{1. Generation of the double-stranded fibre.} Are any two double-stranded reconstructions
of a circular sequence joined by transpositions at repeated $k$-mers, inversions at inverted repeats
of length $k$, and global reverse complementation? This is true within one lattice point (Ukkonen,
Pevzner), true on every case that could be closed, and true on the palindromic class of Chapter~10 by
Kotzig's theorem (Chapter~11). The open part is the passage between lattice points: whether any two box
points with connected support are joined by \emph{realizable} inversions.

\paragraph{2. Odd word length.} Chapter~10's reduction needs palindromic $k$-mers, which exist only for
even $k$. For odd $k$ the quotient of the double cover is a signed graph without fixed vertices
(Chapter~6), and a reconstruction is an Eulerian circuit of a bidirected graph. Is the count still
$\#$P-hard?

\paragraph{3. Approximation.} Does the double-stranded count admit a fully polynomial randomized
approximation scheme? The same question for Eulerian circuits of undirected graphs is open, and
Chapter~10 shows that an answer for genomes would answer it for $4$-regular graphs.

\paragraph{4. Parity of flip distances.} In Chapter~7, reconstructions with the same spectrum are joined
by flip paths of even length only, at $k=11$ and $k=12$ for $\phi$X174. We have no explanation.

\paragraph{5. Beyond two copies.} The sandpile group of a repeat skeleton with two-copy repeats is
presented by an interlace matrix (Chapter~3). For three or more copies no pairwise presentation
exists (Chapter~4). Is there a presentation by data of higher order, and what is it?

\paragraph{What is not on this list.} The specifications of several chapters proposed further
directions that the chapters themselves do not support: Massey products on graphs, whose second
cohomology vanishes; a sheaf formalism for phasing, which leaves its complexity unchanged; and a gauge
theory of chromatin, which has no observable that can be derived. Thermodynamic phase transitions and
anyonic representations of sequences were also suggested, but no chapter defines the objects they would
be about. We list them here only to say that, in the form proposed, they are not open problems of this
book. A problem belongs on the list above only once it can be stated in the terms of a chapter.
